\section{Language and Referencing Guide}

This document uses several tools to ensure consistent and high\-/quality language and references. This section provides guidance on correctly referencing sections, figures, code, and literature, as well as on proper punctuation, quotation, and list formatting practices.

\subsection{Referencing}

Cross\-/referencing is done using the \latexinline{cleveref} package, which automatically includes the type of reference (e.g., figure, algorithm, listing) in the citation. For example, \cref{lst:fibonacci} shows how to calculate Fibonacci numbers while \cref{fig:theway} displays an image. Algorithms and procedures are referenced similarly, such as \cref{alg:expBySquaring} and \cref{pro:mult}. \Cref{fig:theway,fig:subfigures,lst:csharpMath,lst:mathMinted,lst:matlabHelloWorld,tab:budget,tab:standard}. \latexinline{cleveref} does not order these (figures then listings then tables), so you need to reference them in alphabetical order yourself.

The \latexinline{hyperref} package enables clickable links throughout the document, including cross\-/references, citations, and the table of contents. The \latexinline{bookmark} package extends this with improved PDF bookmark generation. The \latexinline{refcheck} package warns about unused labels \--- remove \latexinline{norefs} to display the labels.

References to academic works should be made using \latexinline{\cite} or \latexinline{\textcite}. For example, \cite{xx00} discusses some of the concepts used here while \textcite{xx00} provides an in\-/depth explanation. The difference is that \latexinline{\textcite} integrates the citation naturally into the sentence structure. For the possessive form, \latexinline{\textscite} gives \textscite{xx00} explanation. DOIs are typeset as links with \latexinline{\doi}, as in \doi{10.1000/182}.

When sections are unnumbered, i.e., when \latexinline{secnumdepth} is set to \num{0}, \latexinline{\cref} has no number to print, so \latexinline{\nameref} is commonly used instead to refer to a section by its title: \nameref{sec:math}. The \latexinline{lastpage} package provides the number of the last page, \pageref{LastPage}, for page footers such as \enquote{Page 3 of \pageref{LastPage}}.

\subsection{Hyphenation and Punctuation}

Different dashes are used for different purposes in English typesetting:
\begin{itemize}
  \item \textbf{Hyphen (\-/)} is used to join words, as in \textquote[][.]{well\-/known algorithm}
  \item \textbf{En dash (\-/\-/)} is used for numeric ranges, such as "pages 10\--20" or \textquote[][.]{January\==March}
  \item \textbf{Em dash (\-/\-/\-/)} is used as a parenthetical break \--- like this \--- in a sentence.
\end{itemize}

\subsection{Quotations}

To ensure proper quotation formatting across languages and document styles, the \latexinline{csquotes} package should be used. It provides the \latexinline{\enquote} command, which automatically adapts to the chosen language. For example: \enquote{This is a correctly formatted quote.} A longer quotation can be set as a block:
\blockquote{This is a block quotation. It is useful for longer passages that should be visually separated from the main text, and it adapts to the document language automatically automatically automatically automatically automatically automatically automatically automatically.} And back again.

\subsection{The \latexinline{glossaries-extra} Package}

The \latexinline{glossaries-extra} package is useful for defining glossary entries, acronyms, initialisms, and symbols. Basic glossary and acronym usage:

\begin{itemize}
  \item \latexinline{\gls{}} on first use expands to the full form: \gls{api}.
  \item \latexinline{\glspl{}} gives the plural: \glspl{api}.
  \item Subsequent uses are abbreviated: \gls{api}, \gls{cpu}, \gls{os}, \gls{dr}, \gls{etc}.
  \item \latexinline{\glsxtrshort{}} forces the short form: \glsxtrshort{cpu}.
  \item \latexinline{\newinitialism{}} defines initialisms separately from acronyms, appearing in their own glossary list.
  \item \latexinline{\tgls{}} and \latexinline{\tglspl{}} give the long form in title case, which is useful in headings: \tgls{etc} and \tglspl{os}.
  \item Symbols are referenced in the same way: $E = m\gls{c}\sq$.
\end{itemize}

Symbol entries defined with the \latexinline{symbols} glossary type are listed in the symbols glossary. All glossaries are printed via \latexinline{\printglossaries}, which is called in the main file before the appendix.

\subsection{Text Formatting}

The \latexinline{pifont} package provides dingbats such as \ding{51}, \ding{55}, and \ding{43}, which are handy for checklists. With \latexinline{relsize}, text can be made \textlarger{larger} or \textsmaller{smaller} relative to the current size, and $\mathlarger{\sum}$ enlarges an operator.

Long unbreakable strings, such as hashes, can be broken anywhere with \latexinline{seqsplit}: \seqsplit{e3b0c44298fc1c149afbf4c8996fb92427ae41e4649b934ca495991b7852b855}.

The \latexinline{soul} package offers \hl{highlighting}, \st{strikethrough}, \so{letterspacing}, and \ul{underlining} that allows hyphenation, while \latexinline{ulem} offers \uline{single}, \uuline{double}, \uwave{wavy}, \dashuline{dashed}, and \dotuline{dotted} underlines, as well as \sout{struck through} and \xout{crossed out} text.

Numbers can be written out with \latexinline{fmtcount} (\numberstringnum{42}, \Numberstringnum{7}, \ordinalnum{1}, \ordinalnum{22}, \ordinalnum{103}, and \ordinalstringnum{21}), which falls back to English for unsupported languages such as Norwegian, and \latexinline{\Romannumeral} (\Romannumeral{2026}). Words with diacritics have shorthands that adapt to the language: a \naive approach and \Naive Bayes.

\subsection{Lists}

The \latexinline{enumitem} package provides flexible list customization. The template configures five levels of nested enumeration:

\begin{enumerate}
  \item First level (arabic)
  \begin{enumerate}
    \item Second level (alphabetic)
    \begin{enumerate}
      \item Third level (roman)
      \begin{enumerate}
        \item Fourth level (Alphabetic)
        \begin{enumerate}
          \item Fifth level (Roman)
        \end{enumerate}
      \end{enumerate}
    \end{enumerate}
  \end{enumerate}
\end{enumerate}

For now, only four levels of \latexinline{itemize} are defined:

\begin{itemize}
  \item First level.
  \begin{itemize}
    \item Second level.
    \begin{itemize}
      \item Third level.
      \begin{itemize}
        \item Fourth level.
      \end{itemize}
    \end{itemize}
  \end{itemize}
\end{itemize}

Inline lists are sometimes useful: \begin{enumerate*}
  \item Something
  \item something else
  \item something else entirely
\end{enumerate*}.

\subsection{Miscellaneous}

URLs should be typeset using \latexinline{\url{}} from the \latexinline{xurl} package, which allows line\-/breaking at sensible points: \url{https://ctan.org/pkg/xurl}.

The \latexinline{comment} package allows entire blocks of content to be suppressed without deleting them.
\begin{comment}
  This content is inside a comment environment and will not appear in the output.
  It can span multiple lines and contain arbitrary LaTeX commands.
\end{comment}
The block above is invisible in the compiled output.

The \latexinline{pdfpages} package allows including external PDF files directly into the document using \latexinline{\includepdf[pages=-]{filename.pdf}}, where \latexinline{pages=-} includes all pages. This is useful for appending datasheets, signed forms, or other supplementary documents.
